\documentclass[11pt,letterpaper]{article}

% ---------- Packages ----------
\usepackage[T1]{fontenc}
\usepackage[utf8]{inputenc}
\usepackage{mathptmx}          % Times-like font
\usepackage{courier}           % monospaced font for code names
\usepackage[top=2.2cm,bottom=2.4cm,left=2.6cm,right=2.6cm]{geometry}
\usepackage{graphicx}
\usepackage{float}
\usepackage{enumitem}
\usepackage{xcolor}
\usepackage{fancyhdr}
\usepackage{needspace}
\usepackage{titlesec}
\usepackage{tikz}
\usepackage{eso-pic}
\usepackage[hidelinks]{hyperref}

\setlength{\parindent}{0pt}
\setlength{\parskip}{4pt}
\linespread{1.05}

% ---------- Double page border (on every page) ----------
\AddToShipoutPictureBG{%
  \begin{tikzpicture}[remember picture,overlay]
    \draw[line width=2.2pt]
      ([shift={(0.9cm,-0.9cm)}]current page.north west) rectangle
      ([shift={(-0.9cm,0.9cm)}]current page.south east);
    \draw[line width=0.7pt]
      ([shift={(1.05cm,-1.05cm)}]current page.north west) rectangle
      ([shift={(-1.05cm,1.05cm)}]current page.south east);
  \end{tikzpicture}%
}

% ---------- Header / footer ("N | Page") ----------
\pagestyle{fancy}
\fancyhf{}
\renewcommand{\headrulewidth}{0pt}
\renewcommand{\footrulewidth}{0.4pt}
\renewcommand{\footrule}{\hbox to\headwidth{\color{black!20}\leaders\hrule height \footrulewidth\hfill}}
\fancyfoot[L]{\small\textbf{\thepage}\,|\,\textcolor{black!55}{P\,a\,g\,e}}

% ---------- Section style ----------
\titleformat{\section}{\large\bfseries}{\thesection}{0.7em}{\underline}
\titlespacing*{\section}{0pt}{14pt}{6pt}

% ---------- Question heading helper ----------
\newcommand{\question}[1]{%
  \Needspace{5\baselineskip}%
  \vspace{6pt}\noindent\textbf{#1}\par\nopagebreak}

% ---------- Screenshot helper (no caption) ----------
\newcommand{\shot}[2][]{%
  \begin{center}\includegraphics[#1]{images/#2}\end{center}}

\begin{document}

% =====================================================
%                     COVER PAGE
% =====================================================
\thispagestyle{fancy}
\begin{center}
\vspace*{1.2cm}
\fbox{\includegraphics[width=0.62\textwidth]{images/logo.jpg}}

\vspace{2.8cm}

{\Large\bfseries Mobile Application Development}\\[10pt]
{\Large\bfseries (Theory)}\\[10pt]
{\Large\bfseries ( Assignment : 01)}

\vspace{1.3cm}

{\large\bfseries\underline{Submitted By:}}\\[8pt]
{\large Name: Areeba Shehzadi}\\[6pt]
{\large Class : BSCS -- 5A}\\[6pt]
{\large Reg No : 242201007}

\vspace{1.0cm}

{\large\bfseries\underline{Submitted To:}}\\[8pt]
{\large Sir Uzair}

\vspace{1.0cm}

{\large\bfseries\underline{Submission Date / Day:}}\\[8pt]
{\large 5\textsuperscript{th} October 2026 / Monday}
\end{center}

\newpage

% =====================================================
%                 TITLE + CONTENTS
% =====================================================
\begin{center}
{\LARGE\bfseries Publishing an Android App to Google Play: Versioning,\\
Signing and Release Workflow}
\end{center}

\vspace{2pt}
{\small\tableofcontents}

% =====================================================
%                   1  INTRODUCTION
% =====================================================
\section{Introduction}

Publishing an Android app involves more than writing code. The app needs a proper
version number so that updates can be tracked, it must be digitally signed so that
its author can be verified, and it must be uploaded to the Google Play Console
together with a store listing, a privacy policy and several declarations. This
report explains each of these stages, answers the ten understanding questions, and
documents a practical demonstration in which a basic ``Hello World'' app was given
a version, signed, and exported as an Android App Bundle (AAB).

% =====================================================
%                   2  PART 1
% =====================================================
\section{Part 1: APK Versioning}

\question{Q1. What are \texttt{versionCode} and \texttt{versionName}, and why are they important?}
\texttt{versionCode} is a whole number that Android and Google Play use internally
to decide which build of an app is newer. \texttt{versionName} is a text label,
such as ``1.0'' or ``2.3.1'', that is shown to users and developers. Users never
see the \texttt{versionCode}, but the store depends on it to manage updates, while
the \texttt{versionName} helps people understand which release they have installed.

\question{Q2. What happens if you do not increase the \texttt{versionCode} when updating your app?}
Google Play Console rejects an upload whose \texttt{versionCode} is the same as, or
lower than, a previous one. Even if the file were installed manually, the device
would not treat it as an upgrade of the existing app. Every new release must
therefore carry a higher \texttt{versionCode}.

\question{Q3. Where can you find and edit these values in Android Studio?}
Both values are in the module-level \texttt{build.gradle} file (Gradle Scripts
$\rightarrow$ \texttt{build.gradle} (Module :app)) inside the \texttt{defaultConfig}
block. They can also be viewed and edited through the Project Structure dialog
(Ctrl+Alt+Shift+S). After editing, the project must be synced with Gradle.

% =====================================================
%                   3  PART 2
% =====================================================
\section{Part 2: Generating a Signed Build}

\question{Q4. What is a \texttt{.jks} (keystore) file, and why is it critical?}
A Java KeyStore (\texttt{.jks}) file stores the private key used to sign an app.
The signature proves that the app and all its future updates come from the same
developer. If the key is lost, the developer may not be able to publish updates
under the same app identity, and if it is stolen, someone else could sign apps
pretending to be the developer.

\question{Q5. What is the difference between \texttt{.apk} and \texttt{.aab} files?}
An APK is a complete installable package that can be installed directly on a
device. An AAB is a publishing format that is uploaded to Google Play, which then
generates optimized APKs for each device type, for example matching screen density
and CPU architecture. This usually results in smaller downloads for users. Google
Play requires new apps to be published as AAB files.

\question{Q6. Step-by-step: how to generate a signed release file in Android Studio}
\begin{enumerate}[leftmargin=2em,itemsep=4pt]
  \item Open the project and set \texttt{versionCode} and \texttt{versionName} in
        \texttt{build.gradle}, then sync Gradle.
        \shot[width=0.62\linewidth,height=4cm,keepaspectratio]{sc00.png}

  \item Choose Build $\rightarrow$ Generate Signed App Bundle / APK.
        \shot[width=0.9\linewidth,height=1.4cm,keepaspectratio]{sc01.png}
        \vspace{-8pt}
        \shot[width=0.6\linewidth,height=6.2cm,keepaspectratio]{sc02.png}

  \item Select Android App Bundle and click Next.
        \shot[width=0.8\linewidth,height=6.4cm,keepaspectratio]{sc03.png}

  \item Click Create new\ldots{} next to the key store path.
        \shot[width=0.8\linewidth,height=6.2cm,keepaspectratio]{sc04.png}

  \item Enter the keystore location and file name, a keystore password, a key alias,
        a key password, a validity of 25 years, and at least the developer's name
        in the certificate section. Click OK.
        \shot[width=0.8\linewidth,height=8.2cm,keepaspectratio]{sc05.png}

  \item Click Next, select the release build variant, and click Create/Finish.
        \shot[width=0.8\linewidth,height=6.2cm,keepaspectratio]{sc06.png}

  \item Wait for Gradle to finish. A notification confirms success, and the file
        \texttt{app-release.aab} is created under
        \texttt{app/build/outputs/bundle/release/}.
        \shot[width=\linewidth,height=3.2cm,keepaspectratio]{sc07.png}
\end{enumerate}

\question{Q7. What precautions should you take with your keystore file and passwords?}
The keystore and its passwords should be backed up in at least two safe places,
such as cloud storage and an external drive. They should never be uploaded to a
public repository like GitHub or shared with other people. Passwords should be
stored in a password manager or another secure place. Developers can also enrol in
Play App Signing, where Google holds the final app signing key, so a lost upload
key can be reset.

% =====================================================
%                   4  PART 3
% =====================================================
\section{Part 3: Publishing to Google Play}

\question{Q8. What is required to open a Google Play Developer account?}
A Google account is needed, along with a one-time registration fee of \$25 paid by
card, acceptance of the Developer Distribution Agreement, and identity verification
using a government-issued ID, proof of address and phone verification. Personal
accounts created after 13 November 2023 must also complete a closed test with at
least 12 testers for 14 continuous days before they can apply for production
access. Organization accounts are treated differently and need business details.

\question{Q9. Major steps in publishing an app on the Google Play Console}
\begin{enumerate}[leftmargin=2em,itemsep=2pt]
  \item Create a developer account and complete verification.
  \item Click Create app and enter the app name, language, app or game type, and
        free or paid status.
  \item Complete the dashboard tasks: privacy policy, app access, ads declaration,
        content-rating questionnaire, target audience, and data safety form.
  \item Prepare the store listing: title, short and full description, app icon,
        feature graphic and screenshots.
  \item Create a release on a testing track (internal or closed) and upload the
        signed \texttt{.aab} file.
  \item For new personal accounts, run the closed test with 12 testers for 14 days,
        then apply for production access.
  \item Create a production release, submit it for review, and roll it out once
        approved.
\end{enumerate}

\question{Q10. Common errors or rejections and how to avoid them}
\begin{itemize}[leftmargin=2em,itemsep=2pt]
  \item \textbf{Duplicate or low \texttt{versionCode}:} always increase it for
        every upload.
  \item \textbf{Using \texttt{com.example} as the package name:} choose a unique
        application ID before the first upload.
  \item \textbf{Missing or invalid privacy policy:} host a public policy page and
        add its link.
  \item \textbf{Outdated target API level:} Google requires new apps and updates to
        target a recent Android version (API 36 from 31 August 2026), so check
        \texttt{targetSdk}.
  \item \textbf{Uploading an APK instead of an AAB:} export a signed AAB.
  \item \textbf{Inaccurate declarations or misleading descriptions:} fill in the
        data safety form and content rating honestly and read the Play policies
        before submitting.
  \item \textbf{Not meeting the testing requirement:} recruit enough testers early
        and ask them to actually use the app during the 14 days.
\end{itemize}

% =====================================================
%                   5  PRACTICAL TASK
% =====================================================
\section{Practical Task}

A basic ``Hello World'' application named MyApplication was created in Android
Studio with the application ID \texttt{com.areeba.hellorelease}. The version was set
to \texttt{versionCode} 1 and \texttt{versionName} ``1.0'', with \texttt{minSdk} 24
and \texttt{targetSdk} 37. A signed release bundle was then generated.

\shot[width=\linewidth,height=7cm,keepaspectratio]{sc08.png}

\Needspace{8\baselineskip}
\subsection*{Step 1: Version settings}
Figure~\ref{fig:gradle} shows the \texttt{defaultConfig} block in
\texttt{build.gradle}.

\begin{figure}[H]
  \centering
  \includegraphics[width=0.85\linewidth,height=4.5cm,keepaspectratio]{images/sc09.png}
  \caption{\texttt{versionCode} and \texttt{versionName} in \texttt{build.gradle}}
  \label{fig:gradle}
\end{figure}

\Needspace{24\baselineskip}
\subsection*{Step 2: Generate Signed Bundle}
Figure~\ref{fig:bundle} shows the Generate Signed Bundle window with the release
build variant selected.

\begin{figure}[H]
  \centering
  \includegraphics[width=0.8\linewidth,height=7.2cm,keepaspectratio]{images/sc10.png}
  \caption{Generate Signed App Bundle window}
  \label{fig:bundle}
\end{figure}

\Needspace{20\baselineskip}
\subsection*{Step 3: Resulting AAB file}
Figure~\ref{fig:aab} shows the generated \texttt{app-release.aab} file (about
4.63~MB) in the output folder.

\begin{figure}[H]
  \centering
  \includegraphics[width=\linewidth,height=6.5cm,keepaspectratio]{images/sc11.png}
  \caption{Generated \texttt{app-release.aab} file}
  \label{fig:aab}
\end{figure}

% =====================================================
%                   6  PLAY CONSOLE PLAN
% =====================================================
\section{Play Console Publishing Plan and Current Status}

The signed Android App Bundle (\texttt{app-release.aab}) is ready for upload. The
remaining publishing workflow in the Google Play Console is summarised below.

\begin{itemize}[leftmargin=2em,itemsep=2pt]
  \item \textbf{Ready:} application ID \texttt{com.areeba.hellorelease},
        \texttt{versionCode} 1, \texttt{versionName} ``1.0'', target SDK 37, and a
        signed release AAB.
  \item \textbf{Next:} create a personal developer account (one-time \$25 fee) and
        complete identity verification.
  \item \textbf{Then:} create the app, complete the store listing, privacy policy,
        content rating, target audience and data safety sections.
  \item \textbf{Testing:} upload the AAB to a closed testing track and keep at least
        12 testers opted in for 14 continuous days, as required for new personal
        accounts.
  \item \textbf{Release:} apply for production access, then roll out the release
        after Google's review.
\end{itemize}

Because of the mandatory 14-day closed testing period, a public production release
cannot be completed immediately after the first upload.

% =====================================================
%                   7  CONCLUSION
% =====================================================
\section{Conclusion}

This assignment covered the full path from a finished app to a store-ready release.
Version codes make updates traceable, the keystore proves the developer's identity,
and the AAB format lets Google Play deliver smaller, device-specific installs. The
Play Console adds further requirements such as a privacy policy, declarations,
identity verification and, for new personal accounts, a 14-day closed test.
Understanding and preparing for these steps in advance helps avoid common rejections
and delays.

\end{document}
